\section{The sharp fixed-field moment for the cubic pair}\label{app:cubic-sharp}
The degree-two internal factors require a separate sharp moment,
because Theorem~\ref{thm:config-mixing} assumes degree at least three.
We keep the uniform configuration law throughout and retain every
lattice and Gaussian prefactor.

\begin{theorem}[Cubic raw moment]\label{thm:cubic-raw}
For every fixed odd prime $p$, two independent degree-two internal
configurations and any prescribed cross perfect matching satisfy
\begin{equation}\label{eq:cubic-raw}
 \E_{\rm config}|\ker_{\Fp}(A-I)|=2p+o_p(1)
\end{equation}
along all $N\to\infty$. After conditioning the internal factors
simple,
\[
 \limsup_N\E_{\rm simple}|\ker_{\Fp}(A-I)|
 \le(1+e^{3/2})p.
\]
Both conclusions are uniform over prescribed matchings.
\end{theorem}

\subsection{Exact profiles and the only entropy-zero phases}
Relabel the second block through its cross matching. Independence
and relabeling invariance leave two independent internal configurations,
and $A=\left(\begin{smallmatrix}C_1&I\\I&C_2\end{smallmatrix}\right)$.
After conjugating by $\diag(I,-I)$, the kernel equations are
$C_1x=x+z=C_2z$. At one identified coordinate choose
$A,D,B,E\in\Fp$; the other two ordered internal labels are
$C=A+D-B$ and $F=A+D-E$.
Let $W_1,W_2$ be the local integer table increments
\[
 (W_1)_{ab}=\one_{A=a}(\one_{B=b}+\one_{C=b}),\qquad
 (W_2)_{de}=\one_{D=d}(\one_{E=e}+\one_{F=e}).
\]
For a table $m$ let
$P(m)=\one_{\{m=m^T,\,m_{aa}\text{ even}\}}
\prod_{a<b}m_{ab}!\prod_a(m_{aa}-1)!!$.
For $N$ iid uniform local profiles and $M_j=\sum_iW_{j,i}$,
the exact distinguished-stub count is
\begin{equation}\label{eq:cubic-iid}
 \E|\ker_{\Fp}(A-I)|
 =\frac{p^{4N}}{((2N-1)!!)^2}\E[P(M_1)P(M_2)].
\end{equation}
Every pairing and vector determine exactly one ordered assignment,
including loops with diagonal weight two.

An equivalent two-star count retains the common central-pair counts
$n_{ad}$ and star counts $n^1_{adb},n^2_{ade}$. Its assignment
factor is
$N!\prod_{ad}n_{ad}!/(\prod_{adb}n^1_{adb}!\prod_{ade}n^2_{ade}!)$.
This follows by choosing central labels and then both ordered stars
within each central class. With laws $q_1,q_2$, common marginal
$\pi$, and symmetric edge laws $\rho_1,\rho_2$, each type term has
logarithm $NF+O_p(\log(N+1))$, uniformly at zero entries, where
\[
 F=H(q_1)+H(q_2)-H(\pi)-H(\rho_1)-H(\rho_2).
\]
Symmetrize $q_1$ under $B\mapsto A+D-B$ and $q_2$ under
$E\mapsto A+D-E$, obtaining $s_1,s_2$. The tables and $\pi$ are
unchanged, and entropy gain is $\KL(q_j\Vert s_j)$. The chain rule gives
\begin{equation}\label{eq:cubic-zero-identity}
 F=-I_\pi(A;D)-I_{s_1}(D;B\mid A)-I_{s_2}(A;E\mid D)
                -\KL(q_1\Vert s_1)-\KL(q_2\Vert s_2)\le0.
\end{equation}
At equality the roots are independent and retained neighbors depend
only on their own root. If the law of $D$ has two support points,
the two resulting slot reflections force the first neighbor law
to be invariant under a nonzero field translation, hence uniform.
Symmetry of its edge law forces the law of $A$ to be uniform; the
reverse argument makes the other star uniform as well.
If $D$ is a point at $c$, symmetry of the second edge law forces
both its neighbors to be $c$, and their sum makes $A=c$.
Thus there are exactly $p+1$ zero phases: the uniform two-star law
and the $p$ planted constants. The feasible law space is compact.
Its complement outside fixed neighborhoods of these phases has
negative entropy maximum and contributes exponentially little,
including the polynomial number of types and factorial remainders.

\subsection{Sparse perturbations of planted constants}
Subtract a planted constant and let $k$ count nonzero individual
coordinates. Collapse each cross-matching pair to a supervertex.
If $m$ pairs meet the support and $s$ of them are singletons, then
$m=(k+s)/2$. Use all internal edges incident to the support as a
certificate. If $e$ lie within it and $B$ leave it, its edge count
is $E=e+B=k+B/2\le2k$.
Each singleton pair's zero partner needs a support neighbor to cancel
the nonzero matching contribution. A touched zero coordinate in an
otherwise untouched pair needs two support edges. If there are $u$
such zero coordinates and $r$ new supervertices, then
$B\ge s+2u$ and $r\le u$. Hence, for $v=m+r$,
\begin{equation}\label{eq:cubic-sparse-excess}
 E-v\ge k/2.
\end{equation}
This allows arbitrary field cancellations, loops, and repeated edges.

For $k\le N/4$, choose locations in at most $N^{m+r}/(m!r!)$
ways, occupancy patterns in $3^m$ ways, and nonzero labels in at
most $p^k$ ways. There are at most $4v\le8k$ selected internal
stubs and at most $2^{8k}(4k)^E$ partial pairings. Any valid colored
prescription has probability at most $N^{-E}$, since every remaining
pair denominator is at least $N$. By \eqref{eq:cubic-sparse-excess},
\[
 \frac{(4k)^EN^{m+r-E}}{m!r!}
 \le16^k(k/N)^{k/2}\frac{k^m}{m!}\frac{k^r}{r!}.
\]
Sum $m,r$ using $\sum_jk^j/j!=e^k$ and at most $2k$ choices
of $E$. With $C_*=12288e^2$, the expected count of support-$k$
kernel vectors is at most
$2k(C_*p)^k(k/N)^{k/2}$. Support one is impossible: a zero whole
block forces the other zero through the identity matching.
For sufficiently small fixed $\epsilon_p$, the consecutive ratio
is at most $1/2$ through $2\le k\le\epsilon_pN$.
Their total expected count is $O_p(N^{-1})$.
The $p$ translated neighborhoods are disjoint and contribute
$p+O_p(N^{-1})$. Small fixed central-pair neighborhoods of each
point phase lie inside them.

\subsection{Covariance and the maximal joint support lattice}
Let $\mathcal H$ be real $p\times p$ tables of total zero and
$\Pi$ its row-constant orthogonal projection. Put
\[
 g=p(p-1)/2,\quad s=g+p-1,\quad r_0=p^2-1=g+s.
\]
For a total-zero test table $f$,
\[
 \Var(f(A,B)+f(A,C))
 =\frac2{p^2}\|f\|^2+\frac2{p^3}\sum_a(\sum_bf_{ab})^2.
\]
Conditional on $A$, $B,C$ are independent uniform. Each of
$(A,B),(A,C)$ is independent of each of $(D,E),(D,F)$, since the
corresponding four linear forms are invertible coordinates of
$A,D,B,E$. Thus the two table blocks have zero cross covariance,
and on $\mathcal H$
\begin{equation}\label{eq:cubic-covariance}
 \Sigma=\frac2{p^2}(I+\Pi),\qquad
 \det\Sigma=(4/p^2)^{p-1}(2/p^2)^{2g}.
\end{equation}
In particular the joint covariance has full real rank.

Let $L_0$ be the integer total-zero lattice in $\mathcal H^2$,
of covolume $p^2$, and let $L$ be generated by differences of
local increments $(W_1,W_2)$. Each table has even row sums.
Writing their half row sums as $n^1_a,n^2_d$, all support points obey
the three independent congruences modulo $p$
\begin{align*}
 L_1&=\sum_{ab}b m^1_{ab}-\sum_a a n^1_a-\sum_d d n^2_d=0,\\
 L_2&=\sum_{de}e m^2_{de}-\sum_a a n^1_a-\sum_d d n^2_d=0,\\
 Q_*&=\sum_{ab}ab m^1_{ab}-\sum_{de}de m^2_{de}
                  -\sum_a a^2n^1_a+\sum_d d^2n^2_d=0.
\end{align*}
The row parities have index $2^{2p-2}$. The field constraints are
independent on that lattice: the first-table coefficient of $b$
in a dependence is the $L_1$ coefficient plus $a$ times the $Q_*$
coefficient, forcing both zero for every $a$, then forcing the
$L_2$ coefficient zero. Oddness of $p$ makes row-even variations
surject onto the fixed-total field tables.

There are no further restrictions. A character constant on the
local support has, modulo constants on the full tables,
\begin{align*}
 \theta^1_{ab}&=\alpha_a+(ta+k)b/p,&
 2\alpha_a&=R-[ta^2+(k+l)a]/p,\\
 \theta^2_{de}&=\gamma_d+(-td+l)e/p,&
 2\gamma_d&=S+[td^2-(k+l)d]/p,
\end{align*}
with $t,k,l\in\Fp$, and phases interpreted modulo one.
Indeed varying $B$ at fixed $A,D$ forces each row to be a constant
plus an additive field character; the same holds in the second
table. Rectangular mixed differences in $a,d$ force the slopes
$ta+k,-td+l$, and the remaining equations give the displayed halves.
Conversely every such character is constant on the support.
There are $p^3$ slope choices and $2^{2p-2}$ independent half-sign
choices after removing global table constants. Full rank and lattice
duality therefore give exactly
\begin{equation}\label{eq:cubic-L}
 [L_0:L]=2^{2p-2}p^3,\qquad\operatorname{covol}(L)=p^5 2^{2p-2}.
\end{equation}

\subsection{The valid symmetric lattice and all Gaussian constants}
Let $\mathcal S$ be symmetric real total-zero tables, of dimension
$s$. Its full integer lattice has covolume $p2^{g/2}$: the full
symmetric integer lattice contributes $2^{g/2}$ and the primitive
total-sum functional has Euclidean gradient of norm $p$.
Even diagonals impose index $2^{p-1}$; after imposing them, even
row sums impose another independent index $2^{p-1}$, the rank of
the complete graph's degree-parity map over $\F_2$.
For symmetric tables $L_1=-L_2=T_*$. The constraints $T_*,Q_*$
remain independent: changing first-table diagonal entries by $+2$
at $(a,a)$ and $-2$ at $(0,0)$ changes them by $(a,a^2)$, which
span $\Fp^2$ using $a=1,-1$.
Consequently the valid difference lattice $K$ in $\mathcal S^2$ has
\begin{equation}\label{eq:cubic-K}
 \operatorname{covol}(K)=p^4 2^{g+4p-4},\qquad
 \frac{\operatorname{covol}(L)}{\operatorname{covol}(K)}
       =p2^{-g-2p+2}.
\end{equation}
The affine valid lattice is nonempty at every $N$: put both tables'
only nonzero entry $2N$ at $(0,0)$. Its difference lattice is $K$,
so its central mesh has density $N^s/\operatorname{covol}(K)$,
with no residue or parity subsequence restriction.

Set $t_N=2N/p^2$ and $m=t_N\one\one^T+\sqrt N X$.
Uniformly for symmetric even-diagonal $m$ with $\|X\|\le\log N$,
Stirling gives
\begin{equation}\label{eq:cubic-Stirling}
 \frac{P(m)}{(2N-1)!!}
 =p^{-2N}(4\pi N/p^2)^{g/2}2^{(p-1)/2}
                e^{(p^2/8)\|X\|^2}(1+o_p(1)).
\end{equation}
Off-diagonal factorials provide $(2\pi t_N)^{g/2}$, diagonal
double factorials supply their $\sqrt2$ factors, and the denominator
contributes one inverse $\sqrt2$. The cubic exponent error is
$O_p(\log^3N/\sqrt N)=o(1)$, and the fixed-total linear term vanishes.

For a finite-support iid law with positive covariance and maximal
difference lattice $L$, the relative local limit at
$\|m-N\mu\|\le\sqrt N\log N$ is its Gaussian density times
$\operatorname{covol}(L)(1+o(1))$. Here is the needed justification.
Tilt to mean $m/N$ using the inverse function theorem; the tilt is
$O(\log N/\sqrt N)$. Maximality of $L$ leaves a single maximum on
its reciprocal torus. Outside a fixed neighborhood the modulus is
uniformly below one; near zero positive covariance supplies uniform
Gaussian domination. Fourier inversion at the tilted mean, splitting
at frequency $N^{-2/5}$, gives the exact covolume prefactor with
uniform relative error. Undoing the tilt changes the quadratic
Gaussian exponent by $O(\log^3N/\sqrt N)$. This also proves eventual
positive mass at every central affine-lattice point, without assuming
semigroup normality.

Apply this to the joint increment law. Its mass is
$\operatorname{covol}(L)(2\pi N)^{-r_0}(\det\Sigma)^{-1}$
times the two-block Gaussian and $(1+o(1))$.
Adding the pairing weight \eqref{eq:cubic-Stirling} modifies each
block's precision on $\mathcal S$ to eigenvalues $p^2/8$ with
multiplicity $p-1$ and $p^2/4$ with multiplicity $g$.
Indeed $\Sigma^{-1}=(p^2/2)(I-\Pi/2)$, and the compression of
$\Pi$ to $\mathcal S$ has eigenvalue $1/2$ on
$X_{ab}=v_a+v_b$, $\sum v_a=0$, and zero on symmetric row-sum-zero
tables. The one-block integral is therefore
\[
 I=(2\pi)^{s/2}(8/p^2)^{(p-1)/2}(4/p^2)^{g/2}.
\]
In \eqref{eq:cubic-iid}, the powers of $N$ cancel because
$g+s=r_0$. The remaining coefficient is
\begin{equation}\label{eq:cubic-p-coefficient}
 (4\pi/p^2)^g2^{p-1}(2\pi)^{-r_0}(\det\Sigma)^{-1}
 \frac{\operatorname{covol}(L)}{\operatorname{covol}(K)} I^2=p.
\end{equation}
For example $I^2/\det\Sigma=(2\pi)^s2^{p-1}p^{2g}$, so the
analytic factors are $2^{g+2p-2}$, exactly canceled by
\eqref{eq:cubic-K} except for $p$.

This logarithmic window covers the entire contribution of a sufficiently
small fixed uniform-table neighborhood. The finite-support rate
function is analytic there with Hessian the inverse covariance.
Subtracting the log pairing weight leaves the strictly positive
precision just computed, whose least eigenvalue is $p^2/8$.
Shrink the fixed neighborhood to absorb cubic terms. Chernoff and
uniform Stirling bound each remaining weighted term by
$C_pN^g\exp[-c_p\|m-t_N(\one\one^T,\one\one^T)\|^2/N]$.
Polynomially many tables outside the $\sqrt N\log N$ window thus
contribute $o(1)$. Inside it the affine-lattice Gaussian Riemann
sum converges with the density in \eqref{eq:cubic-K}.
The whole uniform neighborhood contributes $p+o_p(1)$.

\subsection{Region closure and simple conditioning}
Choose the planted central-pair neighborhoods small enough to lie
inside the sparse translated balls, and the uniform edge-table
neighborhood disjoint from them. In the compact feasible two-star
space, the complement contains none of the $p+1$ zeros of
\eqref{eq:cubic-zero-identity}. Its contribution is exponentially
small. The uniform-table sector counts all types with those tables,
not only profiles assumed close to uniform; the sparse sector counts
all vectors near constants regardless of their neighbor profiles.
These sectors therefore cover the same exact moment without an
unexamined intermediate region. Their contributions are $p+o(1)$,
$p+O(N^{-1})$, and $o(1)$, proving \eqref{eq:cubic-raw}.

One degree-two configuration is simple with probability tending to
$e^{-3/4}$. This follows from the fixed-degree simplicity formula,
or from the coefficient ratio of
$(1-z)^{-1/2}e^{-z/2-z^2/4}$ and $(1-z)^{-1/2}$; each simple
graph has $2^N$ distinguished-stub realizations. Both blocks are
simple with probability $s_N\to e^{-3/2}$, exactly giving independent
uniform simple degree-two graphs. Subtracting the deterministic
planted $p$ vectors before conditioning gives
$\E_{\rm simple}K\le p+(\E_{\rm config}K-p)/s_N$, proving the
simple-model bound in Theorem~\ref{thm:cubic-raw}.
Gauge invariance was exact for every prescribed matching, proving
its uniformity assertion as well.
